\section{Planificación en el espacio latente: cómo las predicciones del modelo se convierten en acciones}

El propio modelo del mundo solo responde a la pregunta ``¿qué sucede si realizo esta secuencia de acciones?''; el controller debe buscar a la inversa ``¿qué acción debo tomar para llegar al objetivo?''. LeWorldModel utiliza control predictivo basado en modelos (MPC): realiza rollout repetidos de acciones candidatas en el espacio latente, calcula la distancia entre el punto final y el objetivo, hace que el solver actualice la secuencia de acciones y luego ejecuta solo los primeros pasos antes de volver a observar el entorno.

\subsection{Bucle cerrado de predicción-optimización del Latent MPC}

La diapositiva 45 debe leerse de izquierda a derecha y luego desde abajo derecha hacia las acciones: primero, la imagen actual y la imagen objetivo se mapean al mismo encoder en $z_1, z_g$; luego, el predictor genera $\hat z_2, \ldots, \hat z_H$ basándose en la secuencia de acciones; el costo compara el punto final con el objetivo; finalmente, el solver actualiza las acciones y repite el ciclo. Dado que el encoder no genera píxeles, cada rollout candidato solo se ejecuta sobre un pequeño espacio latente.

\lecturefigure{slide-045.jpg}{Bucle cerrado de control predictivo basado en modelos (model-predictive control) en el espacio latente de LeWorldModel}{CS25 V6 official deck, p.~45}

\[
\hat z_{h+1}=P_\phi(\hat z_h,a_h),\quad \hat z_1=E_\theta(o_1),
\qquad
J(a_{1:H})=\|\hat z_{H+1}-E_\theta(o_g)\|_2^2+\beta R(a_{1:H}).
\]

En esta ecuación, $H$ es el horizonte de planificación, $o_1$ es la observación actual, $o_g$ es la imagen objetivo, $a_{1:H}$ es la secuencia candidata de acciones, $J$ es el costo de planificación, $R$ representa una penalización por suavidad de acción, amplitud o límites, y $\beta$ es el peso de restricción. En la práctica, el MPC real re-codifica la observación tras ejecutar el primer paso para reducir los errores acumulados en bucle abierto (open-loop).

\begin{lstlisting}[language=Python,caption={Ciclo mínimo de predicción-optimización del Latent MPC}]
actions = initialize_action_sequence(horizon)
for _ in range(num_solver_steps):
    state = encoder(current_observation)
    for action in actions:
        state = predictor(state, action)
    cost = latent_distance(state, encoder(goal_observation))
    cost += action_constraint_penalty(actions)
    actions = solver.update(actions, cost)
execute(actions[0])
\end{lstlisting}

\teachervoice{00:42:59--00:47:04, Lucas divide la evaluación en control online y física intuitiva. Al planificar, primero se muestran o inicializan trayectorias de acciones, luego se realizan rollout en el espacio latente y finalmente se actualizan las acciones usando la distancia del goal embedding; esto constituye exactamente la interfaz entre el modelo del mundo y el controller.}

\begin{knowledgebox}{Uso por primera vez: Control Predictivo Basado en Modelos (Model Predictive Control)}
El MPC compromete cada vez solo un horizonte corto: predice el futuro, optimiza las acciones, ejecuta el primer paso, obtiene una nueva observación y vuelve a planificar. La diferencia con un plan abierto (open-loop) de una sola vez es la corrección de errores en bucle cerrado continua; la diferencia con una política directa es que, en tiempo de prueba, aún se llama explícitamente al modelo del mundo y al optimizador, por lo que es más lento pero maneja mejor los objetivos nuevos.
\end{knowledgebox}

\subsection{No basta comparar el éxito promedio en los cuatro entornos}

La diapositiva 46 abarca Two-Room, Reacher, Push-T y OGBench-Cube, con dificultades y geometrías de acción diferentes. LeWM supera a PLDM end-to-end en ciertas tareas complejas y compite con DINO-WM que utiliza DINOv2; sin embargo, no es óptimo en todos los entornos más simples. Se deben examinar por separado el éxito final, la varianza, si se utiliza propiocepción (proprioception) y la carga computacional de cada planificación.

\lecturefigure{slide-046.jpg}{Resultados de control en múltiples entornos: tasas de éxito y diferencias entre tareas de los diferentes baselines}{CS25 V6 official deck, p.~46}

\begin{knowledgebox}{Lectura de la figura: cuatro preguntas para comparar los baselines de control}
En primer lugar, ¿la entrada es solo píxeles o incluye también propiocepción? En segundo lugar, ¿el encoder congela un modelo fundamental (foundation model) o se entrena end-to-end? En tercer lugar, ¿el solver, el horizonte y el número de candidatos son consistentes? En cuarto lugar, ¿las diferencias en la tasa de éxito superan a los márgenes de error? Solo controlando estas condiciones es posible interpretar velocidad y tasa de éxito en un mismo gráfico de Pareto.
\end{knowledgebox}

\subsection{La aceleración de planificación 48 veces proviene del modelo compacto, no es una solución mágica}

La diapositiva 47 coloca el tiempo total de planificación junto a la tasa de éxito. Con aproximadamente 15M de parámetros, un espacio latente compacto y menos tokens, LeWM permite que un ciclo de búsqueda de acciones en una sola GPU dure menos de un segundo, logrando así un aumento de velocidad (speedup) máximo de aproximadamente $48\times$ frente a los modelos del mundo basados en modelos fundamentales. Este número depende de la configuración fija de planificación; cambiar el solver, el tamaño del lote (batch size), el horizonte o el hardware alterará esta proporción.

\lecturefigure{slide-047.jpg}{Velocidad de planificación y tasa de éxito: compensación a nivel de sistema de LeWorldModel}{CS25 V6 official deck, p.~47}

\teachervoice{00:47:04--00:49:08, el docente atribuye la ventaja a una representación end-to-end compacta: los muchos tokens generados por el encoder del modelo fundamental hacen que cada imaginación sea más costosa. Simultáneamente, en clase se reconoce que LeWM es competitivo pero no lidera de manera absoluta en todas las tareas.}

\begin{warningbox}{La velocidad de planificación no puede separarse de la configuración del controller}
Que el ``modelo sea más pequeño'' solo reduce el costo de cada rollout; el retraso total también se multiplica por el número de trayectorias candidatas, las iteraciones del solver, el horizonte y la frecuencia de replaneamiento. Si para compensar errores del modelo se requieren más muestras, los beneficios end-to-end podrían reducirse. Los informes de ingeniería deben proporcionar simultáneamente los FLOPs del modelo, el presupuesto del solver, el tiempo real (wall-clock) y la tasa de éxito.
\end{warningbox}

\subsection{Resumen de la sección}

El Latent MPC convierte el modelo del mundo en un controller ejecutable: el encoder define el estado y el objetivo, el predictor realiza rollout de las consecuencias de las acciones, y el costo junto con el solver buscan a la inversa las acciones. Las principales ventajas sistémicas de LeWM provienen del espacio latente compacto y del modelo pequeño, lo que abarata los numerosos rollout candidatos; su evidencia es un Pareto velocidad-tasa de éxito bajo entornos específicos y un presupuesto fijo del controller, no una garantía incondicional de control en tiempo real.


